\documentclass[12pt]{article}
\usepackage[utf8]{inputenc}
\usepackage[T1]{fontenc}
\usepackage{hyperref}
\title{Yoron Language Resources for K-12 Curriculum Development}
\author{Compiled by Claude Sonnet 5.5 (claude-sonnet-5-5) for Andr\'as Kornai}
\date{}
\begin{document}
\maketitle

\section{Introduction}
Yoron (Yunnu Futuba; ISO 639-3 \texttt{yox}; Glottocode yoro1243) is a Northern Ryukyuan language of Yoronjima, Kagoshima Prefecture, Japan. It is one of the most endangered languages in Japan, with about 950 speakers, all adults, as of 2004. It was long treated as a dialect of Amami or Kunigami, but a growing number of authors treat it as its own language.

\section{Primary Educational Resources}
No school textbooks or primers were found in search results.

\section{Linguistic and Reference Materials}
\begin{itemize}
\item Wikipedia, ``Yoron language'': \url{https://en.wikipedia.org/wiki/Yoron_language}
\item Wikipedia, ``Yoronjima'': \url{https://en.wikipedia.org/wiki/Yoronjima}
\item Omniglot page on Yoron writing. It reportedly links to a PDF on written Yoron by Wolfram Siegel; I did not open it. \url{https://omniglot.com/writing/yoron.htm}
\item JLect database entries: \url{https://www.jlect.com/entry/5/yunnufutuba/} and \url{https://www.jlect.com/entry/386/yunnunchu}
\item Glosbe crowd-sourced translations (thin): \url{https://glosbe.com/ask/yox}
\item Hayashi Yuuka, a paper on the potential form in the Yoron dialect and the grammatical person of the subject (2014): \url{https://cir.nii.ac.jp/crid/1050845763866855040}
\item A repository record found in search (contents not verified): \url{https://iuk-repo.repo.nii.ac.jp/records/647}
\end{itemize}

\section{Religious and Community Resources}
None found. The search summary mentions a local dictionary compiled by Kiku Chiyo (born 1927), which records the Mugiya-higashiku form of the language. I did not find it online.

\section{Digital and Audio Resources}
No audio or video resources were found.

\section{Curriculum Development Context}
Yoron is highly endangered and has no known school programme. Descriptive work exists in Japanese academic literature, and the Kiku Chiyo dictionary is a local lexical source worth tracing. A curriculum would need community and speaker involvement and standardization across the dialect groups (Chabana, Mugiya, and others).
\end{document}